\section*{Hoofdstuk \ref{ch:g10:exercise-and-energy} --- Inspanning en de energiebehoefte van het lichaam}

\begin{solution}{exo:g10:exercise-and-energy:1}
De energie die in volledige rust wordt verbruikt om het lichaam in leven
te houden --- hart, hersenen, temperatuur, onderhoud van de \omterm{def:g10:cells-common-unit:cell}{cellen}:
ongeveer \qty{80}{W}, zo'n \qty{7000}{kJ} per dag.
\end{solution}

\begin{solution}{exo:g10:exercise-and-energy:2}
$2.0 \times 20 = \qty{40}{kJ}$ per minuut, dus $40000/60 \approx
\qty{670}{W}$.
\end{solution}

\begin{solution}{exo:g10:exercise-and-energy:3}
De fosfaatvoorraad van de spier (ongeveer tien seconden); de
melkzuurgisting (één tot twee minuten maximale inspanning); de
\omterm{prop:g10:cell-metabolism:respiration}{celademhaling} (uren).
\end{solution}

\begin{solution}{exo:g10:exercise-and-energy:4}
Het hoogste tempo waarmee het lichaam zuurstof kan opnemen en gebruiken.
$3500/70 = \qty{50}{mL/min}$ per kilogram.
\end{solution}

\begin{solution}{exo:g10:exercise-and-energy:5}
De spieren, ongeveer \qty{400}{g}, als brandstof voor de spier die het
opslaat; de lever, ongeveer \qty{100}{g}, als glucose aan het bloed
afgegeven voor de hersenen en de andere weefsels.
\end{solution}

\begin{solution}{exo:g10:exercise-and-energy:6}
\qty{2.7}{L/min}: $2.7 \times 20 = \qty{54}{kJ}$ per minuut, dus
\qty{900}{W}. Warmte: $900 - 200 = \qty{700}{W}$.
\end{solution}

\begin{solution}{exo:g10:exercise-and-energy:7}
Werk $mgh = 60 \times 9.8 \times 800 \approx \qty{470}{kJ}$. Bij 25\%
rendement verbruikten de spieren ongeveer \qty{1900}{kJ}. Over
\qty{7200}{s}: ongeveer \qty{260}{W} boven de rust.
\end{solution}

\begin{solution}{exo:g10:exercise-and-energy:8}
Een sprint van 100 meter draait vrijwel volledig op de fosfaatvoorraad:
geen zuurstofschuld, weinig melkzuur. Een race van 400 meter duurt
ongeveer een minuut en haalt de helft van haar energie uit de
melkzuurgisting: melkzuur hoopt zich op (de pijn) en de ademhaling die
het moet opruimen houdt de loper daarna nog minuten zwaar ademend.
\end{solution}

\begin{solution}{exo:g10:exercise-and-energy:9}
Koolhydraatenergie $0.6 \times 3000 = \qty{1800}{kJ}$; \omterm{def:g10:exercise-and-energy:glycogen}{glycogeen}
$1800/17 \approx \qty{106}{g}$; vrijgekomen water ongeveer \qty{320}{g}.
\end{solution}

\begin{solution}{exo:g10:exercise-and-energy:10}
Metabool vermogen ongeveer \qty{600}{W}: $\num{1.1e6}/600 \approx
\qty{1830}{s}$, ongeveer 30 minuten.
\end{solution}

\begin{solution}{exo:g10:exercise-and-energy:11}
Drie kwart van de energie die in de spieren vrijkomt is warmte, tijdens
zware inspanning enkele honderden watt. Zweten voert haar af: water
verdampen kost ongeveer \qty{2.4}{kJ} per gram, en de bloedstroom naar
de huid brengt de warmte daarheen. Zonder zweten zou de
lichaamstemperatuur om de paar minuten met een graad stijgen.
\end{solution}

\begin{solution}{exo:g10:exercise-and-energy:12}
Bergop is het benodigde vermogen evenredig met de lichaamsmassa: de
proefpersoon van \qty{60}{kg} heeft \qty{60}{mL/min} per kilogram tegen
\num{40}, en klimt sneller. Op de laboratoriumfiets hangt het vermogen
niet van de massa af, hebben beiden dezelfde \qty{3.6}{L/min}, en kunnen
beiden dezelfde watts volhouden.
\end{solution}

\begin{solution}{exo:g10:exercise-and-energy:13}
De ontbrekende \qty{0.4}{L/min} zuurstofequivalent wordt door de
melkzuurgisting in de spieren geleverd. Melkzuur hoopt zich op; binnen
enkele minuten branden de spieren en moet de loper terugvallen naar een
tempo dat de ademhaling alleen kan dekken.
\end{solution}

\begin{solution}{exo:g10:exercise-and-energy:14}
Korte maximale uitbarstingen draaien op de fosfaatvoorraad en de
\omterm{def:g10:cell-metabolism:fermentation}{gisting}, die helemaal geen vet gebruiken; vet wordt alleen door de
ademhaling verbrand, die minuten nodig heeft om op volle snelheid te
komen en pas bij langdurige matige inspanning overheerst. De totale
energie van enkele uitbarstingen is bovendien klein vergeleken met een
uur gestage inspanning.
\end{solution}

\begin{solution}{exo:g10:exercise-and-energy:15}
Tekort $25000 - 8500 - 6000 = \qty{10500}{kJ}$, dus ongeveer
\qty{280}{g} vet. Eten tijdens de rit houdt de bloedglucose op peil voor
de hersenen en spaart het \omterm{def:g10:exercise-and-energy:glycogen}{glycogeen} voor de beklimmingen, waar het
benodigde vermogen ver boven ligt dan wat vet alleen kan leveren; het vet
dekt het gestage deel van de dag.
\end{solution}

\begin{solution}{pb:g10:exercise-and-energy:1}
\textbf{1.} $4.0 \times 60 \times 42.2 \approx \qty{10100}{kJ}$.

\textbf{2.} $\qty{3}{h}\,\qty{30}{min} = \qty{12600}{s}$:
$\num{1.01e7}/12600 \approx \qty{800}{W}$.

\textbf{3.} $10100/20 \approx \qty{506}{L}$; over \qty{210}{min},
ongeveer \qty{2.4}{L/min}.

\textbf{4.} $2400/60 = \qty{40}{mL/min}$ per kilogram, dus $40/52
\approx 77\%$ van haar maximum.

\textbf{5.} $(2.4 - 0.2)/2.4 \approx 92\%$.

\textbf{6.} $380 \times 17 = \qty{6460}{kJ}$.

\textbf{7.} De wedstrijd kost \qty{240}{kJ} per kilometer: $6460/240
\approx \qty{27}{km}$.

\textbf{8.} Koolhydraatverbruik $0.8 \times 240 = \qty{192}{kJ/km}$:
$6460/192 \approx \qty{34}{km}$ --- de muur.

\textbf{9.} Vetenergie $0.2 \times 10100 \approx \qty{2020}{kJ}$, dus
$2020/38 \approx \qty{53}{g}$: ongeveer 0.6\% van haar \qty{9}{kg}.

\textbf{10.} Vet wordt alleen door de ademhaling verbrand en met een
beperkte snelheid: het kan ruwweg de helft leveren van het vermogen dat
haar tempo vraagt. Zonder \omterm{def:g10:exercise-and-energy:glycogen}{glycogeen} kunnen de spieren niet snel genoeg
ATP maken, en zakt het tempo.

\textbf{11.} $60 \times 3.5 = \qty{210}{g}$, dus $210 \times 17 \approx
\qty{3570}{kJ}$, goed voor $3570/192 \approx \qty{19}{km}$
koolhydraatverbruik.

\textbf{12.} Beschikbaar \omterm{def:g10:chemistry-of-life:families}{koolhydraat} $6460 + 3570 = \qty{10030}{kJ}$:
$10030/192 \approx \qty{52}{km}$, voorbij de finish. Ja.

\textbf{13.} $+\qty{200}{g}$ \omterm{def:g10:exercise-and-energy:glycogen}{glycogeen} $+ \qty{600}{g}$ water:
\qty{0.8}{kg}, wat $0.8 \times 4.0 \approx \qty{3}{kJ}$ extra per
kilometer kost --- verwaarloosbaar tegenover 240.

\textbf{14.} \omterm{def:g10:exercise-and-energy:glycogen}{Glycogeen} $580 \times 17 = \qty{9860}{kJ}$ plus de drank:
ongeveer \qty{13400}{kJ}; koolhydraatverbruik $0.55 \times 240 =
\qty{132}{kJ/km}$: ruim \qty{100}{km}. De muur is ver voorbij de finish
geschoven, met een ruime marge.

\textbf{15.} De hersenen verbranden alleen glucose, geleverd door het
\omterm{def:g10:exercise-and-energy:glycogen}{glycogeen} van de lever. Is die voorraad leeg, dan daalt de bloedglucose,
en de hersenen melden het alarm als duizeligheid en overweldigende
vermoeidheid, wat de spieren ook nog bevatten.

\textbf{16.} $3.6 \times 55 \times 42.2 \approx \qty{8360}{kJ}$;
$8360/20 = \qty{418}{L}$ over \qty{125}{min}: ongeveer \qty{3.3}{L/min}.

\textbf{17.} $3340/55 \approx \qty{61}{mL/min}$ per kilogram, dus
$61/80 \approx 76\%$ --- hetzelfde aandeel als bij Nadia. Het plafond
verschilt, niet het deel ervan dat wordt benut.

\textbf{18.} De loopeconomie zijn de kosten per kilogram per kilometer.
Voor een loper van \qty{55}{kg} scheelt $(4.0 - 3.6) \times 55 \times 42.2 \approx
\qty{930}{kJ}$ over de wedstrijd, ongeveer 9\%.

\textbf{19.} \omterm{def:g10:exercise-and-energy:glycogen}{Glycogeen} \qty{8500}{kJ}; benodigd \omterm{def:g10:chemistry-of-life:families}{koolhydraat} $0.7 \times
8360 \approx \qty{5850}{kJ}$: ja, de voorraden volstaan zonder te
drinken.

\textbf{20.} Zonder voorbereiding raakt het \omterm{def:g10:exercise-and-energy:glycogen}{glycogeen} van Nadia rond
kilometer 34 op; \omterm{def:g10:chemistry-of-life:families}{koolhydraat} stapelen vóór de wedstrijd, \omterm{def:g10:chemistry-of-life:families}{koolhydraat}
drinken tijdens de wedstrijd, en training die het verbrande aandeel vet
verhoogt, duwen de muur elk voorbij de finish.
\end{solution}
